\section{Results}
\label{sec:results}

\subsection{Experimental Setup}





\paragraph{Models.} 

The \textsc{LoopMTP} model is a GPT-2-style decoder-only transformer. In our experiments, we set the number of layers to $L=12$, the embedding dimension to $d=1024$, the number of attention heads to 32, and the nominal FFN hidden dimension to 4096. We employ rotary positional embeddings \citep{su2024roformer}, apply RMSNorm to queries and keys prior to attention \citep{henry2020query, dehghani2023scaling}, and use SwiGLU \citep{shazeer2020glu} in the FFNs. 
We compare \textsc{LoopMTP} against two baselines: (1) a non-looped baseline, and (2) LoopFormer \citep{jeddi2026loopformer}, the current state-of-the-art among looped models. The non-looped baseline keeps the same architectural hyperparameters as \textsc{LoopMTP}, except that its nominal FFN hidden dimension is increased to 4352 (realized width 3072), which more than compensates for the parameters introduced by our gating module and input-fusion projection, giving the baseline roughly 6M more parameters than \textsc{LoopMTP}.
At the same depth and width, LoopFormer has approximately 20M more parameters than our model.


\paragraph{Training.}

All models are trained on the \textit{high-quality} subset of \texttt{Nemotron-CC-v2} \citep{nvidia2025nvidianemotronnano2} and \texttt{Nemotron-CC-Math-v1} \citep{mahabadi2025nemotron} for 6.8B tokens. The models are optimized using the Muon \citep{liu2025muon} and AdamW \citep{loshchilov2017decoupled} optimizers with a peak learning rate of $1.9\times10^{-3}$, a cosine warmup over $0.001$ of total training steps, and a cosine decay schedule down to $1/10$ of the peak learning rate. We train LoopFormer on the same data. Because LoopFormer training was highly unstable across loop configurations, we ran a grid search over learning rate, weight decay, and warmup ratio, yielding 12 configurations per loop count, and report the best result for each.
Appendices~\ref{sec:appendix:implementation_details},~\ref{sec:appendix:WD}, and ~\ref{sec:flops}
provide further details on training and implementation for both approaches, the role of weight decay, and a FLOPs comparison.






\paragraph{Evaluation.}

In our experiments, we report perplexity on a random 5M-token subset of FineWeb-Edu \citep{lozhkov2024fineweb-edu} and OpenWebText \citep{Gokaslan2019OpenWeb}. Downstream performance is evaluated on three groups of tasks using the OLMES framework~\citep{gu2025olmes}: (1) commonsense benchmarks: ARC-Challenge (ARC-C) and ARC-Easy (ARC-E) \citep{allenai:arc}, HellaSwag (HS) \citep{zellers2019hellaswag}, LAMBADA (LB) \citep{paperno2016lambada}, PIQA \citep{bisk2020piqa}, SocialIQA (SIQA) \citep{sap2019social}, and Winogrande (WG) \citep{sakaguchi2021winogrande}; (2) a math suite: Algebra, Counting \& Probability, Geometry, Intermediate Algebra, Number Theory, Prealgebra, and Precalculus; and (3) code benchmarks: CodeX HumanEval \citep{chen2021evaluating} and MBPP \citep{austin2021program}. The reported accuracy and BPB values are averaged over 3 random seeds.
Because our experiments are conducted at a small scale, we use bits-per-byte (BPB) as our primary metric, as it is a reliable proxy for downstream performance at larger scales \citep{gadre2025language, heineman2026signal}. 
For benchmarks where accuracy provides a meaningful signal, we report accuracy as well. In Section~\ref{sec:expert}, we report accuracy on GSM8K~\cite{cobbe2021gsm8k} and BPB on the math suite.








\definecolor{lightblue}{RGB}{220,230,255}
\definecolor{lightgray}{RGB}{245,245,245}


    

        

        



        




\begin{table*}[t]
    \centering
    \resizebox{0.994\linewidth}{!}{
    \begin{tabular}{
        @{}l
        c |
        c c |
        c c c c c c c c |
        c c c c@{}
        }
    \toprule
        & & \multicolumn{2}{c}{\textit{Lang.\ Modeling (PPL $\downarrow$)}} & \multicolumn{8}{c}{\textit{General Tasks (Accuracy $\uparrow$)}} & \multicolumn{4}{c}{\textit{Math/Code/QA (BPB $\downarrow$)}} \\
        \cmidrule(lr){3-4} \cmidrule(lr){5-12} \cmidrule(lr){13-16}
        Model & \# Params & FineWeb-Edu & OpenWebText & ARC-C & ARC-E & HS & WG & SIQA & PIQA & LB & \textit{Avg} & QA & Math & Code & \textit{Avg} \\
        \midrule
        \textbf{Non-looped}          & 266M & 21.08 & 23.56 &  31.60 & 59.90 & 38.48 & 50.96 & 44.85 & 65.98 & 32.16 & 46.28 & 0.9948 & 0.7114 & 0.8624 & 0.8562\\
        
        \midrule
        \textit{\textbf{LoopFormer$^*$ }
        } & \multirow{5}{*}{280M} &&&&&&&&&&&&&&\\
        \quad Loops=3$^\dagger$  & &  21.18 & 23.91        & 32.04 & 58.69 & 38.61 & 50.24 & 44.19 & 66.00 & 30.98 & 45.82 & 1.0212 & 0.7344 & 0.9037 & 0.8864\\
        \quad Loops=5  & &  21.05 & 23.65        & 35.24 & 59.65 & 41.61 & 52.35 & 45.87 & 65.29 & 35.13 & 47.88 & 0.9747 & 0.6898 & 0.8348 & 0.8331\\
        \quad Loops=7  & &  20.66 & 23.20        & 30.01 & 46.49 & 34.37 & 51.43 & 42.39 & 60.52 & 20.84 & 40.86 & 1.5140 & 1.7259 & 2.3012 & 1.8470\\
        \quad Loops=9  & &  20.72 & 23.31        & 32.17 & 49.97 & 36.50 & 50.86 & 42.94 & 60.97 & 23.73 & 42.45 & 1.4555 & 1.6034 & 1.8089 & 1.6226\\

        \midrule
        \rowcolor{lightgray}
        \textit{\textbf{\textsc{LoopMTP}}\small(ours)\normalsize} & \multirow{5}{*}{} &&&&&&&&&&&&&&\\
        \rowcolor{lightgray}
        \quad Loops=3     &   & 19.43 & 21.57 & \underline{35.49} & 62.28 & 42.60 & 50.07 & 46.37 & 67.17 & 35.48 & 48.49 & 0.9541 & 0.6681 & 0.8071 & 0.8098\\
        \rowcolor{lightgray}
        \quad Loops=5 & \multirow{-1}{*}{260M}   & 18.89 & 20.97 & 35.32 & 63.38 & 44.35 & 52.54 & 46.88 & 67.52 & \underline{36.79} & 49.54 & 0.9413 & \underline{0.6584} & \underline{0.7921} & \underline{0.7973}\\
        \rowcolor{lightgray}
        \quad Loops=7     &   & \underline{18.87} & \underline{20.94} & 35.10 & \textbf{64.07} & \textbf{44.78} & \textbf{52.99} & \underline{46.95} & \textbf{67.97} & 36.39 & \underline{49.75} & \underline{0.9369} & \textbf{0.6575} & \textbf{0.7822} & \textbf{0.7922}\\
        \rowcolor{lightgray}
        \quad Loops=9     &  & \textbf{18.86} & \textbf{20.92} & \textbf{36.01} & \underline{63.71} & \underline{44.64} & \underline{52.70} & \textbf{47.78} & \underline{67.86} & \textbf{37.47} & \textbf{50.02} & \textbf{0.9366} & 0.6586 & 0.8003 & 0.7985\\
    \bottomrule
    \end{tabular}
    }
    \vspace{-5pt}
    \caption{\textbf{Main results.} We report perplexity for language modeling tasks, accuracy for general tasks, and BPB for QA, math, and code suites. QA reports BPB on the same seven benchmarks for which accuracy is shown under \textit{General Tasks}.
    $^*$ for LoopFormer \cite{jeddi2026loopformer}, stable training required a separate grid search over learning rate, weight decay, and warmup ratio at each loop count. $^\dagger$ indicates runs whose reported averages are over two random seeds, as one of the three runs diverged.
    }
    \vspace{-10pt}
    \label{tab:main_results}
\end{table*}









































\subsection{Downstream Tasks Performance}
\label{sec:main_results}


Table~\ref{tab:main_results} reports model performance on the perplexity benchmarks and downstream tasks. 
The numbers reported in the main experiment are averaged across three random seeds. Detailed results are presented in Appendix~\ref{sec:appendix:multiseed_results}.

Despite being slightly smaller, \textsc{LoopMTP} outperforms the non-looped baseline on nearly every benchmark: not only on reasoning-based tasks,
but also on perplexity, general tasks, and QA suites. On general tasks, it improves average accuracy by up to 8.08\% over the non-looped baseline (50.02\% vs.\ 46.28\%) and 21.76\% over LoopFormer (49.75\% vs.\ 40.86\%). On QA, math, and code tasks, it reduces BPB by up to 7.5\% and 57\% relative to the two baselines (0.7922 vs.\ 0.8562 and 1.8470) and improves perplexity by 10.5\% on FineWeb-Edu and 11.2\% on OpenWebText.
Additionally, at matched loop counts, \textsc{LoopMTP} outperforms LoopFormer in 27 of 28 cases (4 loop counts $\times$ 7 benchmarks), and also achieves better perplexity and QA/math/code performance, without per-loop tuning of stability-critical hyperparameters (learning rate, weight decay, warmup). Only $\lambda_\text{align}$ (which controls the strength of the auxiliary MTP signal) is tuned, and all swept values train stably.
    
\paragraph{\textsc{LoopMTP} scales stably with loop count.}
While the number of unique parameters stays fixed, which governs a looped model's capacity to store factual information \citep{frey2026dual}, perplexity for \textsc{LoopMTP} continues to improve with additional loops.
This is consistent with \citet{zhu2025scaling}, who find that looped models can be stronger at knowledge manipulation.
Perplexity and general-task accuracy improve monotonically in loop count ($T$), whereas QA, math, and code BPB improve through $T{=}7$ and regress slightly at $T{=}9$. Crucially, \textsc{LoopMTP}'s training remains stable, unlike LoopFormer's: one reason for LoopFormer's relatively worse performance is its unstable training across different loop numbers and even random seeds (see Appendix~\ref{sec:appendix:multiseed_results}). 
We leave the analysis of looping limits and upper bounds to future work.






\subsection{MTP Auxiliary Objective Impact}
\paragraph{Effect of MTP on performance.}
Figure~\ref{fig:mtp_effect_performance} (Top) shows the impact of the MTP auxiliary objective on looped transformer performance across loop counts. On the QA/math/code suites, runs with $\lambda_\text{align} {>} 0$ (\textit{w/ MTP}) show clear improvements over runs with $\lambda_\text{align} {=} 0$ (\textit{w/o MTP}). The effect is most pronounced on math, where step-by-step reasoning is required: without MTP, increasing the loop count causes performance to degrade monotonically. General-task average accuracy is also higher with MTP. Together, these results confirm that MTP is an effective guidance signal for looped models, with negligible overhead. Furthermore, the overall results suggest that combining MTP with looped models is a promising direction: unlike \citet{gloeckle2024better}, where MTP hurt performance at a small scale, here it consistently improves results.


\begin{figure}[t!]
    \centering
    \includegraphics[
        width=0.98\linewidth,
        trim=0cm 0pt 0cm 7pt, %
        clip
    ]{figures/rerendered/LoopMTP_vs_loop_bpb_acc.pdf}
    \vspace{-10pt}
    \caption{\textbf{Impact of MTP auxiliary training signal on performance and per-loop representations.} (Top) Aligning per-loop representations with the MTP signal improves performance, particularly at higher loop counts. (Bottom) With MTP, at most iterations yield a more distinct representation.
    }
    \vspace{-10pt}
    \label{fig:mtp_effect_performance}
\end{figure}



\paragraph{Effect of MTP on hidden representations.}
Figure~\ref{fig:mtp_effect_performance} (Bottom) reports cosine similarities for $T{=}5$ models trained with and without the MTP signal: between the input and output of each iteration $r_i$ (left), and between the output of the first iteration and the outputs of subsequent iterations (right). The right panel shows that, for \textit{w/ MTP} model, the second iteration already maps the input to a markedly distinct subspace, and subsequent iterations continue to build on this trajectory. In contrast, the \textit{w/o MTP} model changes representations only gradually across iterations, with similarities remaining relatively high, particularly across the first three iterations. The left panel shows a consistent pattern: at most iterations, the \textit{w/ MTP} model produces a more distinct representation than its non-MTP counterpart,
increasing expressiveness.
These results confirm the existence of \emph{undifferentiated computation}~\citep{yu2025mesh}, visible in the consistently high cross-iteration similarities of \textit{w/o MTP} model, and show that MTP encourages each loop to learn a more unique transformation.


Figure~\ref{fig:MTP_rank} shows the median ground-truth rank (log scale) at each recurrence iteration $t$, evaluated at its trained offset $u_{i+t}$ over 5M tokens from the OpenWebText dataset. This was obtained by feeding the output of each iteration to the language model head and reading off the logit distribution, for a model trained with MTP (\textit{w/ MTP}) and its counterpart trained without (\textit{w/o MTP}). A model that predicts the ground-truth token perfectly assigns it rank one; the lower the rank, the better the model predicts the ground-truth label. The pattern is clear: the rank of the \textit{w/o MTP} model is up to more than an order of magnitude worse than that of the \textit{w/ MTP} model. The \textit{w/ MTP} model, by contrast, stays within a similar range across all future-token predictions. We attribute this to the absence of a discounting factor in our objective: all future tokens are weighted equally.
This confirms the alignment is actually achieved and that it transfers to the language-model head's coordinate system, rather than being satisfied in a subspace the head ignores.



\begin{figure}[t]
    \centering
    \includegraphics[
        width=0.94\linewidth,
        trim=0cm 0pt 0cm 7pt, %
        clip
    ]{figures/mtp_alignment_compare_aligned_L9_vs_L9_noMTP_trend_median_pct.pdf}
    \vspace{-10pt}
    \caption{\textbf{MTP training sharpens ground-truth retrieval.} Despite the lightweight nature of the latent MTP guidance, the \textit{w/ MTP} model achieves a ground-truth rank up to 35.6$\times$ better than the \textit{w/o MTP} model. 
    }
    \vspace{-10pt}
    \label{fig:MTP_rank}
\end{figure}




\subsection{Gating Mechanism Impact}

\begin{figure*}[t!]
    \centering
    \includegraphics[
        width=0.91\linewidth,
        trim=0cm 0pt 0cm 11pt, %
        clip
    ]{figures/rerendered/gate_ablate_table_grid.pdf}
    \vspace{-11pt}
    \caption{\textbf{Gating mechanism effects.} (Left) When \textit{all gates are learnable} we consistently achieve the best performance, particularly as the number of loops increases.
    (Right) Later iterations exhibit higher and more variable gate values, indicating increasingly input-dependent blending.
    }
    \label{fig:gate_distribution}
\end{figure*}





\paragraph{Gates effectiveness.} Figure~\ref{fig:gate_distribution} (Left) illustrates the importance of the gating mechanism. We compare our per-iteration gating, \textit{All (learnable)}, against three baselines. 
In \textit{Only last}, each iteration overwrites the previous one, with the final gate value fixed to one and all others set to zero. \textit{Only last (learnable)} similarly zeros all non-final gates but allows the final gate value to be determined by a context-dependent learnable mechanism. \textit{All (uniform)} combines the representations by simply taking an average.
\textit{All (learnable)} consistently outperforms all other variants. While general tasks accuracy for \textit{Only last}, \textit{All (uniform)}, and \textit{All (learnable)} is nearly identical, BPB reveals more nuance. Math BPB highlights the importance of representation aggregation, particularly as the number of loops increases: while all other variants either considerably diverge or stagnate, \textit{All (learnable)} continues to decrease BPB with more loops. For QA BPB, the limitation of \textit{Only last} is clear: more iterations means more overwriting and greater information loss. This motivates a mechanism to retain information across loops: even one as simple as averaging, \textit{All (uniform)}, can perform well at higher loop counts. 
This pattern holds for Code as well, where our gating consistently outperforms all variants across loop counts.



\paragraph{Gate behaviour.}
Figure~\ref{fig:gate_distribution} (Right) shows the distribution of normalized gate values for a model with $T{=}5$ across recurrence iterations, evaluated on 5 million randomly sampled tokens from FineWebEdu.
Two trends emerge. First, the mean gate value increases with iteration depth: iterations 1--3 maintain a relatively low and stable mean (${\approx}\,0.17$--$0.19$), whereas iterations 4 and 5 rise to ${\approx}\,0.21$ and ${\approx}\,0.26$, respectively.
This suggests that the gating mechanism learns an \emph{iteration-aware} strategy: early gates act conservatively, yielding a stable scaling value, while later loops integrate newly computed information more dependent on the context.
Second, the spread of the distribution grows markedly in later iterations.
Early iterations exhibit tight, unimodal distributions, indicating that the gate behaves uniformly across tokens.
By contrast, iterations 4 and~5 display substantially wider distributions, implying that the gate adopts a more \emph{input-dependent} policy at greater depth; selectively deciding, on a per-token basis, how much iterative information to preserve.
Taken together, these patterns confirm that the gating mechanism is not merely a static interpolation but learns a structured, depth-dependent blending strategy that enables \textsc{LoopMTP} to benefit from additional recurrence steps without destabilizing training.













\section{Small Domain-Specific Expert}
\label{sec:expert}




In this section, we show that, when trained on a specific domain, \textsc{LoopMTP} can be regarded as a small expert model, well suited to domain-specific reasoning or deployment on memory-constrained hardware. We train models on ${\sim}6.8$B tokens from \texttt{Nemotron-CC-Math-v1} dataset on variable number of loops $T\in \{3,7,9,11,13,15\}$.


\begin{figure}[t!]
    \centering
    \includegraphics[
        width=0.98\columnwidth,
        trim=0cm 0pt 0cm 7pt, %
        clip
    ]{figures/rerendered/BPB_vs_ACC_sensitivity_and_ACC.pdf}
    \vspace{-10pt}
    \caption{\textbf{GSM8K accuracy for expert \textsc{LoopMTP} and BPB–accuracy sensitivity.} (Left) \textsc{LoopMTP} achieves 19.03\% vs.\ 7.05\% at equal parameter count. (Right) The steep slope shows that even small BPB reductions yield meaningful accuracy gains.
    }
    \vspace{-10pt}
    \label{fig:expert_acc_perplexity}
\end{figure}

Figure~\ref{fig:expert_acc_perplexity} (Left) shows GSM8K accuracy (8-shot) for models trained with varying numbers of loops. Despite the small model size ($\sim$260M parameters) and limited pretraining budget, and without any finetuning or instruction tuning, our model reaches 19.03\% accuracy on GSM8K. At matched parameter count, the non-looped baseline achieves only 7.05\%; a relative improvement of approximately 170\% under the same parameter memory budget. 

One reading of this gap comes from mechanistic work showing that mathematical ability in LLMs is distributed across a localized parameter subset, with performance scaling with the fraction of that subset removed rather than hinging on a few individual weights \citep{shomali2026llm}. What matters for a math expert may therefore be how much of the reused capacity is allocated to math-relevant computation, rather than the raw parameter count. This is precisely what looping provides: more effective capacity at fixed parameter memory.

Accuracy does degrade at certain loop counts: the number of \emph{hurt} samples (initially correct, flipped to incorrect by additional loops) eventually outweighs the number of \emph{rescued} samples at those loop counts (we suspect this reflects the model's underlying reasoning characteristics; Appendix~\ref{sec:appen:mechanistic_verification} reports related token-length statistics).
Figure~\ref{fig:expert_acc_perplexity} (Right) plots GSM8K accuracy against BPB on the math suite across varying numbers of loops, together with a fitted line. The slope of $-169$ indicates that even a small reduction in BPB translates to considerable accuracy gain. For example, a 0.01 reduction in BPB would lead to roughly a 1.69\% increase in accuracy.












